\documentclass[10pt]{article}

\begin{document}
\title{Project 2 Report}
\author{User~1 and User~2\\\footnotesize\texttt{user@uio.no}}
\date{}
\maketitle

\begin{abstract}
  This is an example. Use this template as a starting point for writing your report for Project 2. The abstract is usually between 100-400 words, and should summarize the contents of your report. In short, state the problem and the methods you have used, and give your results and conclusion.
\end{abstract}

\section{Introduction}
Use the introduction to describe the problem, the methods you will apply to investigate or solve it, any hypotheses you have, and outline your findings.
How you organize the text is up to you, but keep in mind that you are in a sense telling the story of how your project went (without the gritty details). You want to capture the interest of the reader and then keep them interested enough to read your whole paper!
Why is the problem relevant? How are the results useful? What are the implications of your results?
Imagine that you are trying to ``sell'' your ``product'' (i.e. your project and results).

\section{Related Work}
In this section, you contextualize your project by referring to other research that inspired, is similar to, or is otherwise related to your own.
We don't expect you to give a super detailed overview of all research about adaptation methods and ViTs that exists out there, but you can for example try to find some research on the performance of LoRA in different settings.
Put your references in the \texttt{report.bib} file. There should be an example there already \cite{Test}.

\subsection{Important for Project 2}
\label{subsec:important}
In Project 2, we place more emphasis on your literature search, and how you relate other research to the task at hand. This does not mean you need to include as many references as possible, but that each reference should have a \textit{clear} connection to your methodology and your approach. In other words, we want you to motivate your approach based on existing research, and communicate any link to your own work.

You are also encouraged to find particular research that furthers your findings. For instance, if a relevant paper claims that a particular method converges faster with fewer samples, this could be useful to cite in your work if your experiments show similar properties.

\section{Method}
This section describes your methodology in detail. Try to focus on which approach you decided on, and provide a concise explanation of the model architecture, whether or not you used pretrained weights, which libraries you used, and your approach to conditioning on the scores. Use backreferences to previous sections, particularly to the points addressed in Section~\ref{subsec:important} to elucidate on how you tie relevant research to your methodology.

\section{Experiments and Results}
In this section you outline the experiments you have ran, and present the results.
In general, we encourage the use of tables, graphs, and figures.

\section{Discussion}
Here you discuss the results of your experiments. Did it match your initial hypotheses? Why/why not? What are the implications of the results?

\subsection{Important for Project 2}
There are a few points that are of particular importance for project 2.
\begin{itemize}
  \item \textbf{Selecting your approach:} What was your a priori decision process for selecting your approach? How did you weight the different options? What motivated the choice, and how does it relate to the task?
  \item \textbf{Optimal Preferences:} Exactly how does your approach to conditioning generate images that are preferable for the scores of the focus group member(s)? How confident are you in your generated images?
  \item \textbf{Image quality:} How do you evaluate the images generated by your model? Are they sufficiently similar to the images in the original dataset? Why do you think this is?
  \item \textbf{Approach in Hindsight:} Given what you know now, what are the limitations and benefits to your approach compared to the alternatives? Contrast the options, and perform an analysis.
\end{itemize}

\section{Conclusion}
State your conclusion to tie up the report. This is alse where you can mention the direction of future work based on the results of this one.

\bibliographystyle{IEEEtran}
\bibliography{report}

\end{document}